\documentclass[10pt]{article}
\usepackage[a4paper,margin=0.86in]{geometry}
\renewcommand{\rmdefault}{cmr}
\usepackage{amsmath,amssymb,amsthm,mathtools,amscd,booktabs,enumitem,xcolor,hyperref}
\hypersetup{colorlinks=true,linkcolor=blue!55!black,citecolor=blue!55!black,urlcolor=blue!55!black}
\setlength{\emergencystretch}{3em}
\setlength{\parindent}{0pt}
\setlength{\parskip}{0.52em}
\setlist{leftmargin=1.6em,itemsep=0.22em,topsep=0.22em}
\sloppy

\newtheorem{theorem}{Theorem}[section]
\newtheorem{proposition}[theorem]{Proposition}
\newtheorem{corollary}[theorem]{Corollary}
\newtheorem{lemma}[theorem]{Lemma}
\theoremstyle{definition}
\newtheorem{definition}[theorem]{Definition}
\newtheorem{convention}[theorem]{Convention}
\newtheorem{remark}[theorem]{Remark}

\newcommand{\Zl}{\mathbf Z_\ell}
\newcommand{\Ql}{\mathbf Q_\ell}
\newcommand{\Qlb}{\overline{\mathbf Q}_\ell}
\newcommand{\Dbc}{\mathcal D^b_c}
\newcommand{\Dbcnorm}{\mathcal D^b_{c,\mathrm{norm}}}
\newcommand{\Dctf}{\mathcal D^b_{\mathrm{ctf}}}
\newcommand{\Mod}{\mathsf{Mod}}
\newcommand{\Comp}{\operatorname{Comp}}
\newcommand{\Kar}{\operatorname{Kar}}
\newcommand{\Hom}{\operatorname{Hom}}
\newcommand{\pred}[1]{\par\smallskip\begingroup\footnotesize\raggedright
\emph{Immediate predecessors:} #1.\par\endgroup}

\begin{document}
\begin{center}
{\large\bfseries Compactification Coherence for \(Rf_!\)\par}
\medskip
{\normalsize Extension by zero, refinements, composition, and coefficient change\par}
\end{center}

\begin{abstract}
The exceptional direct image is constructed from finite-torsion sheaves,
proper direct image, and compactifications.  Extension by zero begins with
the closed-support condition and passes through exactness to the Verdier
localization.  Dense refinements give compactification independence,
composition, units, and associativity.  Normalized adic systems,
central localization, and idempotent completion then supply the l-adic
coefficient categories.  Every derived category below is an ordinary
triangulated category.
\end{abstract}

\tableofcontents

\section{Finite coefficients and ordinary derived categories}

\begin{convention}[Geometric setting]
Let \(S\) be a noetherian scheme, let \(\ell\) be invertible on \(S\), and
put \(\Lambda_n=\mathbf Z/\ell^n\mathbf Z\).  Schemes are noetherian and of
finite type over \(S\), and morphisms called compactifiable are separated
and of finite type.  The topology \(\tau\) is either the small etale topology
or the small pro-etale topology.  At finite torsion level the constructible
categories on the two sites are identified by the usual comparison
equivalence.
\end{convention}

\begin{definition}[Module sheaves and the derived category]
For a scheme \(T\), let
\[
 \Mod_\tau(T,\Lambda_n)
\]
be the Grothendieck abelian category of sheaves of
\(\Lambda_n\)-modules on \(T_\tau\).  Let
\(\operatorname{Cpx}(T,\Lambda_n)\) be its category of cochain complexes.
The ordinary derived category is the Verdier localization
\[
 D(T,\Lambda_n)=
 \operatorname{Cpx}(T,\Lambda_n)
 [\{\text{quasi-isomorphisms}\}^{-1}].
\]
Its shift is the cochain shift and its distinguished triangles are those
isomorphic to mapping-cone triangles.
\[
 D^+(T,\Lambda_n)\subset D(T,\Lambda_n)
\]
denotes the full triangulated subcategory of complexes whose cohomology
vanishes in all sufficiently negative degrees.
\end{definition}

\begin{definition}[Derived extension of finite coefficients]
For a homomorphism of finite rings \(A\to B\), let \(B_T\) be the constant
ring sheaf on \(T_\tau\).  Write
\(D(T,A)=D(\Mod_\tau(T,A))\), and similarly for \(B\).  A K-flat
replacement \(P_A(K)\to K\) defines
\[
 K\otimes_A^LB:=P_A(K)\otimes_AB_T.
\]
This gives a functor \(D(T,A)\to D(T,B)\), independent of the replacement up to unique natural
isomorphism.  If \(B\) has finite Tor dimension over \(A\), it preserves
bounded-below and bounded constructible objects.  For
\(A\to B\to C\), fix the canonical associator
\[
 (K\otimes_A^LB)\otimes_B^LC
 \xrightarrow{\sim}K\otimes_A^LC
\]
with its identity and threefold-associativity equations.
\end{definition}

\begin{definition}[Constructible and bounded constructible objects]
A sheaf \(F\in\Mod_\tau(T,\Lambda_n)\) is constructible if \(T\) admits a
finite partition by constructible locally closed subschemes
\(T=\coprod_a T_a\) such that \(F|_{T_a}\) is locally constant with finite
stalks.  Define
\[
 \Dbc(T,\Lambda_n)\subset D(T,\Lambda_n)
\]
to be the full subcategory of complexes \(K\) for which every
\(\mathcal H^i(K)\) is constructible and only finitely many of these
cohomology sheaves are nonzero.  It is a full triangulated subcategory of
\(D^+(T,\Lambda_n)\).
\end{definition}

\begin{definition}[Inverse image]
For \(h:T\to T'\), inverse image of module sheaves is the exact functor
\[
 h^*:\Mod_\tau(T',\Lambda_n)\longrightarrow
 \Mod_\tau(T,\Lambda_n).
\]
Applying it termwise gives an exact triangulated functor, again denoted
\(h^*:D(T',\Lambda_n)\to D(T,\Lambda_n)\).  It preserves bounded
below complexes and bounded constructible complexes, and has coherent isomorphisms
\((kh)^*\simeq h^*k^*\) and \(\operatorname{id}^*\simeq\operatorname{id}\).
\end{definition}

\section{Open extension by zero}

\begin{definition}[Support of a section]
Let \(F\) be a sheaf on \(U_\tau\), let \(W\to U\) be an object, and let
\(s\in F(W)\).  The zero locus of \(s\) is the largest open subobject of
\(W\) on which \(s\) restricts to zero.  Its complement is the support
\(\operatorname{supp}(s)\), regarded as a closed subset of \(|W|\).
\end{definition}

\begin{definition}[The support-closed extension presheaf]
Let \(j:U\hookrightarrow T\) be a quasi-compact open immersion and let
\(F\in\Mod_\tau(U,\Lambda_n)\).  For an object \(V\to T\), put
\(V_U=V\times_TU\) and define
\[
 (j^{\mathrm{pre}}_!F)(V)=
 \{s\in F(V_U):\operatorname{supp}(s)\text{ is closed in }V\}.
\]
Restriction is induced by pullback of sections.  A morphism \(F\to G\)
sends a support-closed section to a support-closed section and therefore
defines a morphism \(j^{\mathrm{pre}}_!F\to j^{\mathrm{pre}}_!G\).
\end{definition}

\begin{lemma}[Restriction preserves the support condition]
The preceding restriction maps are well defined, preserve identities and
composition, and make \(j^{\mathrm{pre}}_!F\) a presheaf of
\(\Lambda_n\)-modules on \(T_\tau\).
\end{lemma}
\pred{the support of a section and the support-closed extension presheaf}

\begin{definition}[Sheaf extension by zero]
The extension by zero of \(F\) is the module sheaf
\[
 j_!F:=a_T(j^{\mathrm{pre}}_!F),
\]
where \(a_T\) denotes sheafification in \(\Lambda_n\)-modules.  The
functoriality of the presheaf construction defines an additive functor
\[
 j_!: \Mod_\tau(U,\Lambda_n)\longrightarrow
       \Mod_\tau(T,\Lambda_n).
\]
\end{definition}

\begin{definition}[Unit and counit for the open immersion]
Restriction of support-closed sections defines a natural morphism
\(\eta_F:F\to j^*j_!F\).  For \(G\in\Mod_\tau(T,\Lambda_n)\), extension of
a section of \(j^*G\) when its support is closed defines
\(\epsilon_G:j_!j^*G\to G\).  These maps are natural in \(F\) and \(G\).
\end{definition}

\begin{proposition}[Extension-by-zero adjunction]
For \(F\in\Mod_\tau(U,\Lambda_n)\) and
\(G\in\Mod_\tau(T,\Lambda_n)\), the maps above induce a natural bijection
\[
 \Hom_T(j_!F,G)\xrightarrow{\sim}\Hom_U(F,j^*G).
\]
The two triangle identities hold; hence \(j_!\dashv j^*\).  Moreover,
\(\eta_F:F\to j^*j_!F\) is an isomorphism.
\end{proposition}
\pred{the sheaf extension, its unit, and its counit}

\begin{proposition}[Open and closed restrictions]
Let \(i:Z=T\setminus U\hookrightarrow T\) be the complementary closed
immersion.  There are natural isomorphisms
\[
 j^*j_!\simeq\operatorname{id},\qquad i^*j_!\simeq0.
\]
For a geometric point \(\bar t\to T\),
\[
 (j_!F)_{\bar t}\simeq
 \begin{cases}
 F_{\bar t},&\bar t\text{ factors through }U,\\
 0,&\bar t\text{ factors through }Z.
 \end{cases}
\]
\end{proposition}
\pred{extension-by-zero adjunction and the geometric-point functor}

\begin{lemma}[Additivity and kernels]
The functor \(j_!\) preserves zero objects, finite biproducts, and kernels.
For every morphism \(F\to G\), the canonical map
\[
 j_!(\ker(F\to G))\longrightarrow
 \ker(j_!F\to j_!G)
\]
is an isomorphism.
\end{lemma}
\pred{the stalk formula and exactness of sheafification in module sheaves}

\begin{theorem}[Exactness of extension by zero]
For every short exact sequence on \(U_\tau\),
\[
 0\longrightarrow F'\longrightarrow F\longrightarrow F''\longrightarrow0,
\]
the extended sequence
\[
 0\longrightarrow j_!F'\longrightarrow j_!F
 \longrightarrow j_!F''\longrightarrow0
\]
is exact on \(T_\tau\).  Thus \(j_!\) is an exact additive functor.
\end{theorem}
\pred{additivity and kernels, the stalk formula, and conservativity of geometric stalks}

\begin{proposition}[Finite constructibility under open extension]
If \(F\) is a constructible finite \(\Lambda_n\)-module sheaf on \(U\),
then \(j_!F\) is constructible on \(T\).
\end{proposition}
\pred{the open and closed restrictions and a finite constructible stratification of \(F\)}

\begin{definition}[Termwise extension of complexes]
For a complex \(K=(K^r,d^r)\), define
\[
 J_jK=(j_!K^r,j_!d^r).
\]
For a cochain map or cochain homotopy, apply \(j_!\) in every degree.
This defines an additive functor on complexes and on their homotopy
categories.
\end{definition}

\begin{proposition}[Cohomology and quasi-isomorphisms]
For every complex \(K\) and every integer \(r\), exactness gives a natural
isomorphism
\[
 \mathcal H^r(J_jK)\simeq j_!\mathcal H^r(K).
\]
Consequently \(J_j\) preserves quasi-isomorphisms, boundedness, and bounded
constructible cohomology.
\end{proposition}
\pred{exactness of \(j_!\) and finite constructibility under open extension}

\begin{definition}[Derived extension by zero]
By the universal property of Verdier localization, \(J_j\) factors uniquely
up to unique natural isomorphism as an exact triangulated functor
\[
 Rj_!=j_!:D(U,\Lambda_n)\longrightarrow D(T,\Lambda_n).
\]
Its restriction is an exact functor
\[
 j_!:D^+(U,\Lambda_n)\longrightarrow D^+(T,\Lambda_n),
\qquad
 j_!: \Dbc(U,\Lambda_n)\longrightarrow\Dbc(T,\Lambda_n).
\]
This derived notation is the termwise extension followed by localization;
no resolution enters its definition.
\end{definition}
\pred{termwise extension and preservation of quasi-isomorphisms}

\begin{theorem}[Composition and base change for open extension]
For quasi-compact open immersions \(U\xrightarrow{j}T\xrightarrow{k}V\),
there is a canonical isomorphism
\((kj)_!\simeq k_!j_!\), compatible with threefold composition and identity
open immersions.  In a Cartesian square of open immersions
\[
\begin{CD}
 U'@>{j'}>>T'\\
 @V{g'}VV @VV{g}V\\
 U@>{j}>>T,
\end{CD}
\]
the support pullback map is a natural isomorphism
\[
 \beta_{j,g}:g^*j_!\xrightarrow{\sim}j'_!g'^*.
\]
These isomorphisms satisfy the identity, horizontal-pasting, and
vertical-pasting equations.
\end{theorem}
\pred{the support-closed presheaf, derived extension by zero, and exact inverse image}

\section{Proper direct image input}

\begin{definition}[Derived direct image]
For \(p:T\to Y\), direct image of module sheaves is the left exact functor
\[
 p_*:\Mod_\tau(T,\Lambda_n)\longrightarrow
      \Mod_\tau(Y,\Lambda_n).
\]
Choose a functorial K-injective resolution \(K\to I(K)\) and define
\[
 Rp_*K:=p_*I(K),\qquad
 Rp_*:D(T,\Lambda_n)\longrightarrow D(Y,\Lambda_n).
\]
The right-derived functor preserves bounded-below complexes and therefore
restricts to \(Rp_*:D^+(T,\Lambda_n)\to D^+(Y,\Lambda_n)\).
For composable maps \(T\xrightarrow{p}Y\xrightarrow{r}Z\), derived
adjunction supplies a natural composition isomorphism
\[
 \mu_{r,p}:Rr_*Rp_*\xrightarrow{\sim}R(rp)_*.
\]
\end{definition}

\begin{theorem}[Coherence of proper direct image]
The maps \(\mu_{r,p}\) are natural, equal the identity when one map is an
identity, and satisfy the pentagon equation for three composable maps.
Thus \(T\mapsto D(T,\Lambda_n)\), \(p\mapsto Rp_*\) is an ordinary
triangulated pseudofunctor.
\end{theorem}
\pred{derived direct image and uniqueness of derived functors}

\begin{definition}[Proper base-change morphism]
For a Cartesian square
\[
\begin{CD}
 T'@>{a}>>T\\
 @V{p'}VV @VV{p}V\\
 Y'@>{b}>>Y
\end{CD}
\]
with \(p\) proper, the proper base-change morphism is the mate
\[
 \operatorname{BC}_{p,b}:b^*Rp_*\longrightarrow Rp'_*a^*
\]
obtained from the unit \(\operatorname{id}\to Ra_*a^*\), the Cartesian
identification \(pa=bp'\), the composition maps for direct image, and the
counit \(b^*Rb_*\to\operatorname{id}\).
\end{definition}

\begin{theorem}[Finite-torsion proper base change]
For \(K\in D^+(T,\Lambda_n)\), the morphism
\(\operatorname{BC}_{p,b}(K)\) is an isomorphism.  It is natural in \(K\),
compatible with morphisms of Cartesian squares, and satisfies identity,
horizontal-pasting, and vertical-pasting equations.  For a geometric point
\(\bar y\to Y\), it gives
\[
 (Rp_*K)_{\bar y}\simeq
 R\Gamma(T_{\bar y},K|_{T_{\bar y}}).
\]
\end{theorem}
\pred{the proper base-change morphism, the strict-henselian stalk comparison, and invariance of finite proper-fiber cohomology}

\begin{theorem}[Open-proper interchange]
In a Cartesian square
\[
\begin{CD}
 T'@>{j'}>>\overline T'\\
 @V{p'}VV @VV{p}V\\
 T@>{j}>>\overline T
\end{CD}
\]
with \(j,j'\) quasi-compact open immersions and \(p,p'\) proper, the mate of
proper base change is a natural isomorphism
\[
 \operatorname{Ex}_{j,p}:j_!Rp'_*\xrightarrow{\sim}Rp_*j'_!
\]
on \(D^+(T',\Lambda_n)\).  It is compatible with identities and with
horizontal and vertical pasting of such squares.
\end{theorem}
\pred{open-extension base change, finite-torsion proper base change, and the two adjunctions \(j_!\dashv j^*\), \(p^*\dashv Rp_*\)}

\section{Compactifications and refinements}

\begin{definition}[Compactification]
For \(f:X\to Y\), a compactification is a factorization
\[
 c=(\overline X_c,j_c,p_c,\alpha_c),\qquad
 X\xrightarrow{j_c}\overline X_c\xrightarrow{p_c}Y,
\]
where \(j_c\) is a quasi-compact open immersion, \(p_c\) is proper, and
\(\alpha_c:p_cj_c=f\) is the specified equality.  It is schematically
dense if \(\mathcal O_{\overline X_c}\to(j_c)_*\mathcal O_X\) is injective.
Write \(\Comp_Y(f)\) for the category of schematically dense
compactifications.
\end{definition}

\begin{theorem}[Existence and dense replacement]
Every separated finite-type morphism \(f:X\to Y\) admits a compactification.
For any compactification, the scheme-theoretic closure of \(X\) in its
proper middle term is a schematically dense compactification and maps
properly to the original one over \(Y\).
\end{theorem}
\pred{Nagata compactification and scheme-theoretic closure}

\begin{definition}[Refinement]
For \(c_1,c_2\in\Comp_Y(f)\), a refinement
\(q:c_1\to c_2\) is a proper \(Y\)-morphism
\[
 q:\overline X_{c_1}\longrightarrow\overline X_{c_2}
\]
such that \(qj_{c_1}=j_{c_2}\) and \(p_{c_2}q=p_{c_1}\).  Thus the source
of a refinement dominates its target.
\end{definition}

\begin{definition}[Pullback over the common open]
For a refinement \(q:c_1\to c_2\), define
\[
 P_q:=\overline X_{c_1}\times_{\overline X_{c_2}}X.
\]
Write
\[
 \iota_q:P_q\longrightarrow\overline X_{c_1},\qquad
 \pi_q:P_q\longrightarrow X
\]
for its projections.  The equality \(qj_{c_1}=j_{c_2}\) defines a unique
map
\[
 s_q:X\longrightarrow P_q,\qquad
 \iota_qs_q=j_{c_1},\quad \pi_qs_q=\operatorname{id}_X.
\]
\end{definition}

\begin{lemma}[The projection from the pullback is proper]
The morphism \(\pi_q:P_q\to X\) is proper and separated.
\end{lemma}
\pred{the pullback \(P_q\) and stability of proper morphisms under base change}

\begin{lemma}[The section is an open immersion]
The map \(s_q:X\to P_q\) is an open immersion.  Its image is the open
subscheme \(j_{c_1}(X)\) inside the open subscheme
\(P_q=q^{-1}(j_{c_2}(X))\subset\overline X_{c_1}\).
\end{lemma}
\pred{the pullback \(P_q\), the equality \(qj_{c_1}=j_{c_2}\), and openness of \(j_{c_1}\)}

\begin{lemma}[The section is a closed immersion]
The map \(s_q\) is a closed immersion.
\end{lemma}
\pred{properness and separatedness of \(\pi_q\), and the equation \(\pi_qs_q=\operatorname{id}_X\)}

\begin{lemma}[Density of the section]
The open immersion \(s_q:X\to P_q\) is schematically dense.  Equivalently,
\[
 \mathcal O_{P_q}\longrightarrow (s_q)_*\mathcal O_X
\]
is injective.
\end{lemma}
\pred{schematic density of \(j_{c_1}\) and the fact that \(P_q\) is an open subscheme of \(\overline X_{c_1}\)}

\begin{proposition}[The pullback is the common open]
The morphism \(s_q:X\to P_q\) is an isomorphism.
\end{proposition}
\pred{the open, closed, and schematically dense properties of \(s_q\)}

\begin{corollary}[The source square of a dense refinement]
For a refinement \(q:c_1\to c_2\), the diagram
\[
\begin{CD}
 X@>{j_{c_1}}>>\overline X_{c_1}\\
 @V{\operatorname{id}}VV @VV{q}V\\
 X@>{j_{c_2}}>>\overline X_{c_2}
\end{CD}
\]
is Cartesian.  In particular, open-proper interchange gives
\[
 \xi_q:j_{c_2,!}\xrightarrow{\sim}Rq_*j_{c_1,!}.
\]
\end{corollary}
\pred{the isomorphism \(s_q:X\simeq P_q\) and open-proper interchange}

\begin{definition}[Candidate attached to one compactification]
For \(c\in\Comp_Y(f)\), define
\[
 Rf_{!,c}:=Rp_{c,*}j_{c,!}:
 D^+(X,\Lambda_n)\longrightarrow D^+(Y,\Lambda_n).
\]
The displayed source and target use only exact open extension and
bounded-below derived direct image; no constructibility theorem enters this
definition.
\end{definition}

\begin{definition}[The refinement transformation]
Let \(q:c_1\to c_2\).  Put
\(\gamma_q:=\xi_q^{-1}:Rq_*j_{c_1,!}\to j_{c_2,!}\).  Define
\[
 \begin{aligned}
 \theta_q:Rf_{!,c_1}
 &=Rp_{c_1,*}j_{c_1,!}\\
 &\xrightarrow{\mu_{p_{c_2},q}^{-1}}
 Rp_{c_2,*}Rq_*j_{c_1,!}
 \xrightarrow{Rp_{c_2,*}\gamma_q}
 Rp_{c_2,*}j_{c_2,!}=Rf_{!,c_2}.
 \end{aligned}
\]
\end{definition}

\begin{theorem}[Refinement isomorphism and functoriality]
Every \(\theta_q\) is a natural isomorphism.  For the identity refinement,
\(\theta_{\operatorname{id}}=\operatorname{id}\).  For refinements
\(c_1\xrightarrow{q}c_2\xrightarrow{r}c_3\),
\[
 \theta_{rq}=\theta_r\circ\theta_q.
\]
\end{theorem}
\pred{the source square, the refinement transformation, and coherence of proper direct image and open-proper interchange}

\section{Common refinements and independence}

\begin{definition}[Product-closure refinement]
For \(c_1,c_2\in\Comp_Y(f)\), form
\(\overline X_{c_1}\times_Y\overline X_{c_2}\) and the map
\((j_{c_1},j_{c_2}):X\to
\overline X_{c_1}\times_Y\overline X_{c_2}\).  Let
\(\overline X_{12}\) be the scheme-theoretic closure of its image.  The
restricted projections are denoted
\(r_i:\overline X_{12}\to\overline X_{c_i}\).
\end{definition}

\begin{theorem}[Common dense refinement]
The map \(X\to\overline X_{12}\) is a quasi-compact schematically dense open
immersion, \(\overline X_{12}\to Y\) is proper, and each \(r_i\) is a
refinement.  Hence any finite family of compactifications has a common dense
refinement.
\end{theorem}
\pred{product-closure refinement, stability of proper morphisms under products and closed immersions, and the closure calculation on the common open}

\begin{lemma}[Uniqueness of refinement arrows]
Two refinement arrows having the same source and target are equal.  More
generally, two maps from a schematically dense scheme to a separated target
that agree on the dense open \(X\) are equal.
\end{lemma}
\pred{schematic density and separatedness of the target}

\begin{corollary}[Cofiltered compactification category]
The category \(\Comp_Y(f)\) is nonempty and cofiltered: it has common
refinements, and parallel refinement arrows are equalized.
\end{corollary}
\pred{existence, common dense refinements, and uniqueness of refinement arrows}

\begin{definition}[Comparison through a common refinement]
Let \(d\) dominate \(c_1,c_2\) by refinements \(r_i:d\to c_i\).  Define
\[
 \kappa^{d}_{c_1,c_2}:=
 \theta_{r_2}\circ\theta_{r_1}^{-1}:
 Rf_{!,c_1}\xrightarrow{\sim}Rf_{!,c_2}.
\]
\end{definition}

\begin{theorem}[Independence of the common refinement]
The map \(\kappa^{d}_{c_1,c_2}\) is independent of \(d\) and of its two
refinement arrows.  Denote the resulting natural isomorphism by
\(\kappa_{c_1,c_2}\).  For all \(c_1,c_2,c_3\),
\[
 \kappa_{c,c}=\operatorname{id},\qquad
 \kappa_{c_2,c_1}=\kappa_{c_1,c_2}^{-1},\qquad
 \kappa_{c_2,c_3}\kappa_{c_1,c_2}=\kappa_{c_1,c_3}.
\]
\end{theorem}
\pred{cofilteredness and identity and composition for the transformations \(\theta_q\)}

\begin{definition}[The chosen exceptional direct image]
Choose one compactification \(c_f\in\Comp_Y(f)\) and define
\[
 Rf_!:=Rf_{!,c_f}:
 D^+(X,\Lambda_n)\longrightarrow D^+(Y,\Lambda_n).
\]
For a second choice \(c'_f\), the choice-change isomorphism is
\(\kappa_{c_f,c'_f}:Rf_{!,c_f}\to Rf_{!,c'_f}\).
\end{definition}

\begin{theorem}[Compactification independence]
The choice-change maps are natural, satisfy the identity and cocycle
equations, and are uniquely compatible with all refinement maps.  Thus the
notation \(Rf_!\) is independent of the selected compactification up to one
specified coherent natural isomorphism.
\end{theorem}
\pred{independence and the cocycle for \(\kappa\)}

\section{Composition, identity, and normalization}

\begin{definition}[Successive compactification datum]
Let \(X\xrightarrow{f}Y\xrightarrow{g}Z\), and choose compactifications
\[
 X\xrightarrow{j_f}\overline X_f\xrightarrow{p_f}Y,
 \qquad
 Y\xrightarrow{j_g}\overline Y_g\xrightarrow{p_g}Z.
\]
A successive compactification datum is a Cartesian diagram
\[
\begin{CD}
 \overline X_f@>{a}>>W\\
 @V{p_f}VV @VV{b}V\\
 Y@>{j_g}>>\overline Y_g
\end{CD}
\]
in which \(a\) is a quasi-compact schematically dense open immersion and
\(b\) is proper.  It induces the compactification
\[
 d_{g,f}=\bigl(X\xrightarrow{aj_f}W
 \xrightarrow{p_gb}Z\bigr)
\]
of \(gf\).
\end{definition}

\begin{theorem}[Existence of successive compactifications]
Every pair of chosen compactifications of \(f\) and \(g\) admits a
successive compactification datum.  It may be obtained by compactifying the
map \(j_gp_f:\overline X_f\to\overline Y_g\) and taking the closure whose
restriction over the open \(Y\) is the proper scheme \(\overline X_f\).
\end{theorem}
\pred{Nagata compactification, properness of \(p_f\), and dense replacement}

\begin{definition}[Raw compositor]
For a successive compactification datum, open-proper interchange and the
composition maps for proper direct image define
\[
 \begin{aligned}
 \widetilde c_{g,f}:Rg_{!,c_g}Rf_{!,c_f}
 &=Rp_{g,*}j_{g,!}Rp_{f,*}j_{f,!}\\
 &\xrightarrow{\operatorname{Ex}_{j_g,b}}
 Rp_{g,*}Rb_*a_!j_{f,!}\\
 &\xrightarrow{\sim}
 R(p_gb)_*(aj_f)_!=R(gf)_{!,d_{g,f}}.
 \end{aligned}
\]
\end{definition}

\begin{definition}[Canonical compositor]
For the selected compactification \(c_{gf}\) used to define \(R(gf)_!\), set
\[
 c_{g,f}:=
 \kappa_{d_{g,f},c_{gf}}\circ\widetilde c_{g,f}:
 Rg_!Rf_!\xrightarrow{\sim}R(gf)_!.
\]
\end{definition}

\begin{theorem}[Choice independence of the compositor]
The natural isomorphism \(c_{g,f}\) is unchanged after replacing any chosen
compactification or successive compactification datum, once the canonical
choice-change maps \(\kappa\) are inserted.
\end{theorem}
\pred{raw compositor, compactification independence, and pasting coherence for open-proper interchange}

\begin{definition}[Identity unit]
For \(\operatorname{id}_X\), use the compactification
\(X\xrightarrow{\operatorname{id}}X\xrightarrow{\operatorname{id}}X\).
The two identity functors define
\[
 u_X:R(\operatorname{id}_X)_!\xrightarrow{\sim}
 \operatorname{id}_{D^+(X,\Lambda_n)}.
\]
\end{definition}

\begin{theorem}[Unit and associativity equations]
For every compactifiable \(f:X\to Y\), the two composites involving
\(c_{f,\operatorname{id}}\), \(c_{\operatorname{id},f}\), and the units
\(u_X,u_Y\) are identities.  For
\(X\xrightarrow{f}Y\xrightarrow{g}Z\xrightarrow{h}W\),
\[
 c_{h,gf}\circ(Rh_!c_{g,f})
 =c_{hg,f}\circ(c_{h,g}Rf_!)
\]
as natural transformations \(Rh_!Rg_!Rf_!\to R(hgf)_!\).  Hence
\(f\mapsto Rf_!\) is an ordinary triangulated pseudofunctor.
\end{theorem}
\pred{choice independence of the compositor, the cocycle for \(\kappa\), and associativity of proper direct image and open-proper interchange}

\begin{theorem}[Proper and open normalizations]
If \(f\) is proper, the compactification \((\operatorname{id}_X,f)\) gives
a canonical isomorphism \(Rf_!\simeq Rf_*\).  If \(j:U\hookrightarrow X\)
is a quasi-compact open immersion, the compactification
\((j,\operatorname{id}_X)\) gives \(Rj_!\simeq j_!\).  Both normalizations
commute with the units and compositors above.
\end{theorem}
\pred{the chosen exceptional direct image, exact open extension, and compactification independence}

\section{Constructible landing and cohomological bounds}

\begin{definition}[Compact-support cohomological amplitude]
For \(f:X\to Y\), an interval \([a,b]\) is a compact-support amplitude on
a full subcategory \(\mathcal A\subset\Mod_\tau(X,\Lambda_n)\) if
\(R^qf_!F=0\) for every \(F\in\mathcal A\) and every
\(q\notin[a,b]\).
\end{definition}

\begin{theorem}[Proper direct-image finiteness]
If \(p:T\to Y\) is proper and of finite type, then the ambient functor
restricts to
\[
 Rp_*:\Dbc(T,\Lambda_n)\longrightarrow\Dbc(Y,\Lambda_n).
\]
If the geometric fibers have dimension at most \(d\), then
\(R^qp_*F=0\) for every constructible sheaf \(F\) and every \(q>2d\).
\end{theorem}
\pred{finite-torsion proper base change, finiteness of proper geometric-fiber cohomology, and the proper cohomological-dimension theorem}

\begin{theorem}[Constructibility for one sheaf]
For every separated finite-type \(f:X\to Y\) and every constructible
\(\Lambda_n\)-module sheaf \(F\), each \(R^qf_!F\) is constructible and
all but finitely many are zero.
\end{theorem}
\pred{exact open extension, proper direct-image finiteness, and the ambient compactification formula}

\begin{theorem}[Compact-support finite cohomological dimension]
If every geometric fiber of \(f:X\to Y\) has dimension at most \(d\), then
for every constructible sheaf \(F\),
\[
 R^qf_!F=0\qquad(q>2d).
\]
Equivalently, \([0,2d]\) is a compact-support amplitude on constructible
sheaves.  The bound is stable under arbitrary base change and does not
depend on the dimension of a chosen proper middle term.
\end{theorem}
\pred{the quasi-finite seed, relative affine-line devissage, local factorization, and finite-torsion compact-support base change}

\begin{proposition}[Bounded hypercohomology transfer]
For \(K\in\Dbc(X,\Lambda_n)\), the standard truncation filtration gives a
finite convergent spectral sequence
\[
 E_2^{r,s}=R^rf_!\mathcal H^s(K)
 \Longrightarrow \mathcal H^{r+s}(Rf_!K).
\]
If \(\mathcal H^s(K)=0\) outside \([a,b]\) and the fibers have dimension at
most \(d\), then \(\mathcal H^i(Rf_!K)=0\) outside \([a,b+2d]\), and every
abutment sheaf is constructible.
\end{proposition}
\pred{constructibility for one sheaf, compact-support finite cohomological dimension, and finite truncation triangles}

\begin{corollary}[Finite-level constructible landing]
For every separated finite-type \(f:X\to Y\), the ambient functor restricts
to an exact triangulated functor
\[
 Rf_!: \Dbc(X,\Lambda_n)\longrightarrow\Dbc(Y,\Lambda_n).
\]
When the fibers have dimension at most \(d\), it sends cohomological
interval \([a,b]\) into \([a,b+2d]\).
\end{corollary}
\pred{bounded hypercohomology transfer and the already constructed ambient exceptional direct image}

\begin{corollary}[Finite-level functor package]
At every level \(n\), the functors \(h^*\), \(Rh_*\) for proper \(h\), and
\(Rf_!\) for compactifiable \(f\) act on bounded constructible
triangulated categories.  Their comparison, composition, unit, and
base-change isomorphisms are morphisms inside those categories rather than
only inside the ambient derived categories.
\end{corollary}
\pred{finite-level constructible landing and proper direct-image finiteness}

\section{Integral adic coefficients}

\begin{convention}[Continuous integral coefficients on the pro-etale site]
Let \(E/\Ql\) be finite, let \(\mathcal O_E\) be its ring of integers, let
\(\varpi\) be a uniformizer, and put
\(\Lambda_m=\mathcal O_E/\varpi^m\).  From this section onward the site is
\(T_{\mathrm{pro\acute et}}\).  If \((\Lambda_m)_T\) denotes the finite
constant sheaf of rings on this site, define
\[
 (\mathcal O_E)_T:=\varprojlim_m(\Lambda_m)_T
\]
in sheaves of rings and write \(\mathcal O_E\) for this continuous ring
sheaf.  Let
\(\Mod_{\mathrm{pro\acute et}}(T,\mathcal O_E)\) be the Grothendieck abelian category of
sheaves of modules over \((\mathcal O_E)_T\), and put
\[
 D(T,\mathcal O_E)=D\bigl(\Mod_{\mathrm{pro\acute et}}(T,\mathcal O_E)\bigr).
\]
Every restriction-of-scalars map, tensor product, and completeness
condition below is taken over this ring object.  Derived tensor products
are computed by K-flat replacements in its module-sheaf category.
\end{convention}

\begin{definition}[Strict coefficient systems]
The abelian category
\(\Mod_{\mathrm{pro\acute et}}(T,\mathcal O_{E,\bullet})\) has objects
\[
 M_\bullet=(M_m,\rho_{m+1,m})_{m\ge1},
\]
where \(M_m\) is a sheaf of \(\Lambda_m\)-modules and
\(\rho_{m+1,m}:M_{m+1}\to M_m\) is \(\mathcal O_E\)-linear after restriction
of scalars.  Kernels and cokernels are levelwise.  Its ordinary derived
category is
\[
 D(T,\mathcal O_{E,\bullet})
 :=D\bigl(\Mod_{\mathrm{pro\acute et}}(T,\mathcal O_{E,\bullet})\bigr),
\]
and exact evaluation gives
\(\operatorname{ev}_m:D(T,\mathcal O_{E,\bullet})\to D(T,\Lambda_m)\).
\end{definition}

\begin{definition}[Normalized bounded constructible system]
An object \(K_\bullet\in D(T,\mathcal O_{E,\bullet})\) is normalized if the
transition maps induce coherent isomorphisms
\[
 \operatorname{ev}_{m+1}(K_\bullet)
 \otimes^L_{\Lambda_{m+1}}\Lambda_m
 \xrightarrow{\sim}\operatorname{ev}_m(K_\bullet).
\]
It is bounded constructible with uniform bounds if every level belongs to
\(\Dbc(T,\Lambda_m)\) and one finite cohomological interval works for all
levels.
\end{definition}

\begin{theorem}[Normalized shifts and cones]
Normalized systems with uniform bounded constructible cohomology are closed
under shifts.  If \(u:K_\bullet\to L_\bullet\) is a morphism between two such
systems, then every mapping cone of \(u\) is normalized, its level \(m\) is
isomorphic to
\(\operatorname{Cone}(\operatorname{ev}_m u)\), and one uniform
cohomological interval works for all levels.  Hence these systems form a
full triangulated subcategory
\[
 \Dbcnorm(T,\mathcal O_{E,\bullet})
 \subset D(T,\mathcal O_{E,\bullet}).
\]
\end{theorem}
\pred{exact evaluation, the derived tensor associator, and the two-out-of-three property in a morphism of distinguished triangles}

\begin{definition}[Towers in one coefficient category]
Let
\[
 \mathsf{Tow}_T(\mathcal O_E)=
 \operatorname{Fun}(\mathbf N_{\ge1}^{\mathrm{op}},
                    \Mod_{\mathrm{pro\acute et}}(T,\mathcal O_E)).
\]
Restriction of scalars along
\(\mathcal O_E\to\Lambda_m\) defines an exact functor
\[
 \operatorname{Res}_T:
 \Mod_{\mathrm{pro\acute et}}(T,\mathcal O_{E,\bullet})
 \longrightarrow\mathsf{Tow}_T(\mathcal O_E),
\]
and hence an exact functor on ordinary derived categories.
\end{definition}

\begin{definition}[Derived inverse limit and realization]
The inverse-limit functor
\[
 \varprojlim:\mathsf{Tow}_T(\mathcal O_E)
 \longrightarrow\Mod_{\mathrm{pro\acute et}}(T,\mathcal O_E)
\]
is left exact.  A K-injective resolution in the single abelian tower
category defines
\[
 R\!\varprojlim:
 D(\mathsf{Tow}_T(\mathcal O_E))
 \longrightarrow D(T,\mathcal O_E).
\]
Set
\[
 \operatorname{Real}_T=
 R\!\varprojlim\circ\operatorname{Res}_T.
\]
For each \(m\), the projection of a tower and derived reduction define
\[
 r_m(K_\bullet):
 \operatorname{Real}_T(K_\bullet)
 \otimes^L_{\mathcal O_E}\Lambda_m
 \longrightarrow\operatorname{ev}_m(K_\bullet).
\]
\end{definition}

\begin{theorem}[Adic reconstruction]
For a normalized bounded constructible system \(K_\bullet\), every
\(r_m(K_\bullet)\) is an isomorphism.  The realized object is derived
\(\varpi\)-complete, and a morphism between two realized normalized systems
is an isomorphism if and only if all its finite reductions are
isomorphisms.
\end{theorem}
\pred{the Mittag--Leffler property of the quotient tower, derived completion, and the realization maps \(r_m\)}

\begin{definition}[Integral constructible category]
The category \(\Dbc(T,\mathcal O_E)\) is the full subcategory of
\(D(T,\mathcal O_E)\) consisting of derived \(\varpi\)-complete objects
\(K\) such that
\[
 \operatorname{red}_mK:=
 K\otimes^L_{\mathcal O_E}\Lambda_m
\]
forms a normalized bounded constructible system with one uniform bound.
Adic reconstruction identifies this category with the realized normalized
systems and supplies its ordinary triangulated structure.
\end{definition}

\begin{theorem}[Levelwise exceptional direct image]
Finite-level exceptional direct image and its reduction isomorphisms define
an exact functor
\[
 Rf_{!,\bullet}:
 \Dbcnorm(X,\mathcal O_{E,\bullet})
 \longrightarrow
 \Dbcnorm(Y,\mathcal O_{E,\bullet}),
 \qquad (K_m)_m\longmapsto(Rf_!K_m)_m.
\]
All refinement maps, compactification comparisons, compositors, and units
commute with the transition maps and with finite reduction.
\end{theorem}
\pred{finite-level constructible landing, normalized cones, and derived tensor compatibility for open and proper direct image}

\begin{theorem}[Integral adic exceptional direct image]
Realization defines an exact functor
\[
 Rf_!: \Dbc(X,\mathcal O_E)\longrightarrow\Dbc(Y,\mathcal O_E)
\]
characterized by
\[
 \operatorname{red}_m(Rf_!K)\xrightarrow{\sim}
 Rf_!(\operatorname{red}_mK).
\]
The finite-level refinement, composition, unit, and associativity
isomorphisms lift uniquely to this functor.
\end{theorem}
\pred{levelwise exceptional direct image, adic reconstruction, and conservativity of finite reductions}

\section{Rational and algebraically closed coefficients}

\begin{definition}[Central localization]
Let
\[
 \mathcal A_E(T)=
 \Dbc(T,\mathcal O_E)[\varpi^{-1}]
\]
be the additive category with the same objects as
\(\Dbc(T,\mathcal O_E)\) and
\[
 \Hom_{\mathcal A_E(T)}(K,L)=
 \Hom_{\Dbc(T,\mathcal O_E)}(K,L)\otimes_{\mathcal O_E}E.
\]
The shift is induced from the integral category.  A triangle is
distinguished if it is isomorphic to the image of an integral
distinguished triangle after central localization.
\end{definition}

\begin{theorem}[Triangulation of the central localization]
The preceding shift and distinguished triangles make
\(\mathcal A_E(T)\) an \(E\)-linear triangulated category, and the
localization functor
\(\Dbc(T,\mathcal O_E)\to\mathcal A_E(T)\) is exact.
\end{theorem}
\pred{centrality of multiplication by \(\varpi\) and triangulated localization by a central multiplicative system}

\begin{definition}[Karoubi envelope]
An object of \(\Kar(\mathcal A_E(T))\) is a pair \((K,e)\) with
\(e^2=e\in\operatorname{End}_{\mathcal A_E(T)}(K)\).  A morphism
\((K,e)\to(L,d)\) is a morphism \(u:K\to L\) satisfying
\(u=due\).  Shifts are \((K,e)[1]=(K[1],e[1])\).
\end{definition}

\begin{theorem}[Triangulation of the Karoubi envelope]
The idempotent completion \(\Kar(\mathcal A_E(T))\) has a unique
triangulated structure for which the inclusion of \(\mathcal A_E(T)\) is
exact and every distinguished triangle is a direct summand of a
distinguished triangle from \(\mathcal A_E(T)\).  Define
\[
 \Dbc(T,E):=\Kar(\mathcal A_E(T)).
\]
\end{theorem}
\pred{the triangulated idempotent-completion theorem}

\begin{theorem}[Transport through localization and idempotents]
An exact \(\mathcal O_E\)-linear functor \(F\) on integral constructible
categories extends first to the central localizations and then to the
Karoubi envelopes by
\[
 (K,e)\longmapsto(FK,F(e)).
\]
A natural isomorphism \(\alpha:F\to G\) extends on a summand by
\[
 \alpha_{(K,e)}=G(e)\,\alpha_K\,F(e).
\]
These operations preserve identities, composition, distinguished
triangles, and equations between natural transformations.
\end{theorem}
\pred{triangulation of the central localization and Karoubi envelope}

\begin{definition}[Finite scalar extension]
For a finite extension \(E'/E\), define on an integral representative
\[
 K\otimes_EE':=
 K\otimes^L_{(\mathcal O_E)_T}(\mathcal O_{E'})_T,
\]
and on an idempotent object set
\[
 (K,e)\otimes_EE'=(K\otimes_EE',e\otimes1).
\]
This gives a functor \(\Dbc(T,E)\to\Dbc(T,E')\).
\end{definition}

\begin{theorem}[Exactness and coherence of finite scalar extension]
Finite scalar extension preserves derived completeness, bounded
constructibility, distinguished triangles, and idempotent summands.  Its
identity and tower associators are coherent for every chain of finite
extensions.
\end{theorem}
\pred{finite flatness of rings of integers, finite-reduction compatibility, and transport through localization and idempotents}

\begin{corollary}[Rational exceptional direct image]
The integral functor extends to an exact \(E\)-linear functor
\[
 Rf_{!,E}: \Dbc(X,E)\longrightarrow\Dbc(Y,E),
\]
which we also denote by \(Rf_!\).  Its refinement, composition, unit, and
associativity isomorphisms are obtained by idempotent transport.  For every
finite extension \(E'/E\), it has a
coherent natural isomorphism
\[
 (Rf_{!,E}K)\otimes_EE'\xrightarrow{\sim}
 Rf_{!,E'}(K\otimes_EE'),
\]
compatible with towers of finite extensions and with every displayed
coherence map.
\end{corollary}
\pred{integral adic exceptional direct image and transport through localization and idempotents}

\begin{definition}[Finite coefficient-field pair category]
Fix \(\Qlb\).  Let \(\mathcal C_c(T,\Qlb)\) have objects
\((E,K_E)\), where \(E\subset\Qlb\) is finite over \(\Ql\) and
\(K_E\in\Dbc(T,E)\), and morphisms
\[
 \Hom((E,K_E),(E',L_{E'}))=
 \varinjlim_{M\supset E,E'}
 \Hom_{\Dbc(T,M)}
 (K_E\otimes_EM,L_{E'}\otimes_{E'}M).
\]
The transition maps are exact scalar-extension functors.  Identities and
composition are represented after extension to one common finite coefficient field
and are independent of a further finite extension.
\end{definition}

\begin{definition}[Common-stage shifts, cones, and triangles]
Set \((E,K_E)[1]=(E,K_E[1])\).  If a morphism is represented over a
finite coefficient field \(M\) by
\(u_M:K_E\otimes_EM\to L_{E'}\otimes_{E'}M\), define its common-stage
cone to be \((M,\operatorname{Cone}(u_M))\).  A triangle is distinguished
if it is represented by a distinguished triangle in \(\Dbc(T,M)\) after
extension to one finite coefficient field containing all fields and representatives.
\end{definition}

\begin{theorem}[Independence of common-stage cones]
Two representatives of one morphism acquire isomorphic cones after a
further finite extension.  Their isomorphism class is compatible with
shifts, composition, and scalar extension, and distinguishedness is
independent of the chosen common field.
\end{theorem}
\pred{filteredness of finite subextensions, exact scalar extension, and compatibility of mapping cones with exact functors}

\begin{theorem}[Triangulation of the pair category]
The pair category \(\mathcal C_c(T,\Qlb)\) is a
\(\Qlb\)-linear triangulated category.  Extension from one finite coefficient field to
another is exact, and the definition of a distinguished triangle is
independent of the common finite coefficient field.
\end{theorem}
\pred{independence of common-stage cones and the triangulated axioms after passage to one finite coefficient field}

\begin{definition}[Pair-level exceptional direct image]
Write
\[
 \sigma^!_{M/E,K}:(Rf_{!,E}K)\otimes_EM
 \xrightarrow{\sim}Rf_{!,M}(K\otimes_EM)
\]
for rational scalar compatibility.  Define
\[
 Rf_{!,\mathrm{pair}}(E,K_E)=(E,Rf_{!,E}K_E).
\]
If \([M,u_M]\) represents a morphism
\((E,K_E)\to(E',L_{E'})\), define its image to be represented over \(M\)
by
\[
 (\sigma^!_{M/E',L_{E'}})^{-1}\,
 Rf_{!,M}(u_M)\,\sigma^!_{M/E,K_E}.
\]
\end{definition}

\begin{theorem}[Well-definedness of pair-level exceptional image]
The common-stage formula is independent of the representative and finite
field, preserves identities and composition, and defines an exact
triangulated functor
\[
 Rf_{!,\mathrm{pair}}:
 \mathcal C_c(X,\Qlb)\longrightarrow\mathcal C_c(Y,\Qlb).
\]
The refinement maps, compactification comparisons, compositors, and units
form natural isomorphisms of pair-level functors and retain all their
coherence equations.
\end{theorem}
\pred{scalar compatibility of rational exceptional direct image, independence of common-stage cones, and coherence after one finite extension}

\begin{definition}[Algebraically closed coefficient category]
Define
\[
 \Dbc(T,\Qlb):=\Kar(\mathcal C_c(T,\Qlb)).
\]
Thus an object is a pair \((A,e)\) with \(A\in\mathcal C_c(T,\Qlb)\) and
\(e^2=e\); morphisms, shifts, and triangles are those of the triangulated
Karoubi envelope.
\end{definition}

\begin{definition}[Algebraically closed exceptional direct image]
On an idempotent object \((A,e)\) of the Karoubi envelope, define
\[
 Rf_!(A,e)=
 (Rf_{!,\mathrm{pair}}A,Rf_{!,\mathrm{pair}}(e)).
\]
This gives a functor
\(Rf_!: \Dbc(X,\Qlb)\to\Dbc(Y,\Qlb)\).
\end{definition}

\begin{theorem}[Exactness and coherence over
\texorpdfstring{\(\Qlb\)}{Qbar-l}]
The algebraically closed exceptional direct image is exact.  Refinement
maps, compactification comparisons, compositors, units, and all coherence
equations extend by idempotent transport and are independent of every
common field used to represent the idempotents.
\end{theorem}
\pred{well-defined pair-level exceptional image and transport to the triangulated Karoubi envelope}

\section{Bibliographic remarks}

The geometric ingredients are Nagata compactification, schematic closure,
proper base change, and the product-closure construction of common
refinements.  The coefficient constructions use derived completion,
triangulated central localization, and triangulated idempotent completion.
These results are treated in the Stacks Project chapters on morphisms,
etale and pro-etale cohomology, derived completion, and triangulated
categories.

\end{document}
